\documentclass[11pt]{article}
\usepackage{amssymb}
\usepackage{amsthm}
\usepackage{enumitem}
\usepackage{amsmath, physics}
\usepackage{bm}
\usepackage{adjustbox}
\usepackage{mathrsfs}
\usepackage{graphicx}
% \usepackage{siunitx}
\usepackage{physunits}
\usepackage[mathscr]{euscript}

\title{\textbf{Solved selected problems of Classical Electrodynamics - Hans Ohanian}}
\author{Franco Zacco}
\date{}

\addtolength{\topmargin}{-3cm}
\addtolength{\textheight}{3cm}

\newcommand{\Nat}{\mathbb{N}}
\newcommand{\Z}{\mathbb{Z}}
\newcommand{\Q}{\mathbb{Q}}
\newcommand{\R}{\mathbb{R}}
\newcommand{\diam}{\text{diam}}
\newcommand{\cl}{\text{cl}}
\newcommand{\bdry}{\text{bdry}}
\newcommand{\inter}{\text{int}}
\newcommand{\hatx}{\bm{\hat{x}}}
\newcommand{\haty}{\bm{\hat{y}}}
\newcommand{\hatz}{\bm{\hat{z}}}
\newcommand{\hatr}{\bm{\hat{r}}}
\newcommand{\hatn}{\bm{\hat{n}}}
\newcommand{\hatrho}{\bm{\hat{\rho}}}
\newcommand{\hatphi}{\bm{\hat{\phi}}}
\newcommand{\hattheta}{\bm{\hat{\theta}}}
\newcommand{\vecx}{\bm{x}}
\newcommand{\varep}{\varepsilon}


\theoremstyle{definition}
\newtheorem*{solution*}{Solution}
\renewcommand*{\proofname}{\bf{Solution}}

\begin{document}
\maketitle
\thispagestyle{empty}

\section*{Chapter 5 - Electric Energy}

\subsection*{Exercises}

\begin{proof}{\textbf{Exercise 1.}}
We want to compute the follwoing integral
\begin{align*}
    -\frac{qq'}{4\pi}\frac{1}{|\bm{x'}|}\int \nabla\cdot\nabla\frac{1}{|\bm{x}|}~dV
\end{align*}
So applying Gauss' theorem we have that
\begin{align*}
    -\frac{qq'}{4\pi}\frac{1}{|\bm{x'}|}\int \nabla\cdot\nabla\frac{1}{|\bm{x}|}~dV
    &= -\frac{qq'}{4\pi}\frac{1}{|\bm{x'}|}
    \int \nabla\frac{1}{|\bm{x}|}\cdot \hatn~dS\\
    &= -\frac{qq'}{4\pi}\frac{1}{|\bm{x'}|}
    \int_0^{2\pi}\int_0^{\pi} \pdv{r}(\frac{1}{r})
    ~r^2\sin\theta d\theta d\phi\\
    &= \frac{qq'}{4\pi}\frac{1}{|\bm{x'}|}
    \int_0^{2\pi}\int_0^{\pi} \frac{1}{r^2}~r^2\sin\theta d\theta d\phi\\
    &= \frac{qq'}{4\pi}\frac{1}{|\bm{x'}|}
    [-\cos\pi + \cos 0]~[2\pi - 0]\\
    &= \frac{qq'}{|\bm{x'}|}
\end{align*}
Where in the second equality we used that the normal direction in spherical
coordinates is the radial direction and hence the gradient component we need
is the radial one.
\\
Therefore
\begin{align*}
    U = \text{constant} + \frac{qq'}{|\bm{x'}|}
\end{align*}
\end{proof}

\cleardoublepage
\begin{proof}{\textbf{Exercise 2.}}
Let a charged particle in a circular orbit in a Coulomb field.
\\
Since the particle is on a Coulomb field then the force exerted to it is
\begin{align*}
    \bm{F} = q\bm{E} = q\frac{q'}{r^2}\hatr
\end{align*}
Where $q'$ is the charge that's generating the field and $r$ is the distance
that separates them.
\\
Then the potential energy is given by
\begin{align*}
    - \dv{V}{r} &= F\\
    V &= -qq'\int \frac{1}{r^2}~dr\\
    V &= \frac{qq'}{r}
\end{align*}
On the other hand, since the particle is in circular motion with respect to
the other charge $q'$ then the centripetal force it feels is
\begin{align*}
    \bm{F}_c = -m \frac{v^2}{r}\hatr
\end{align*}
But then, must be that
\begin{align*}
    \bm{F} = q\frac{q'}{r^2}\hatr &= -m \frac{v^2}{r}\hatr = \bm{F}_c
\end{align*}
And hence the tangential velocity $v$ is
\begin{align*}
    v^2 = -\frac{qq'}{mr}
\end{align*}
So the kinetic energy gives us
\begin{align*}
    T = \frac{1}{2}mv^2 = -\frac{1}{2}m\frac{qq'}{mr}
    = -\frac{1}{2}\frac{qq'}{r}
\end{align*}
Since $r$ doesn't change with time because we are in a circular orbit then
the instantaneous kinetic energy is equal to the average kinetic energy and
the instantaneous potential energy is equal to the average potential energy.
\\
Therefore
\begin{align*}
    \overline{T} = -\frac{1}{2}\frac{qq'}{r} = -\frac{1}{2}\overline{V}
\end{align*}
\end{proof}

\cleardoublepage
\begin{proof}{\textbf{Exercise 3.}}
Let us take a Gaussian surface shaped as a box where one face is inside the
metallic plate with charge $Q$ and the other face is inside the dielectric,
then by Gauss law for dielectrics we have that
\begin{align*}
    \int \bm{D}\cdot d\bm{S} &= 4\pi Q_{enc}\\
    \int_{plate} \bm{D}\cdot d\bm{S} + \int_{dielec} \bm{D}\cdot d\bm{S} &= 4\pi Q\\
    0 + D A &= 4\pi Q
\end{align*}
Where we used that $D$ inside the metallic plate is 0 and by symmetry we
neglected the integrals with respect to the other faces of the box.
Therefore
\begin{align*}
    D &= \frac{4\pi Q}{A} = 4\pi \sigma
\end{align*}

\end{proof}

\cleardoublepage
\begin{proof}{\textbf{Exercise 4.}}
Equation (34) states that
\begin{align*}
    F = -\frac{2\pi Q^2}{\epsilon A}
\end{align*}
And equation (40) states that
\begin{align*}
    F = -\frac{\epsilon A}{8\pi}\bigg(\frac{\Delta \Phi}{l}\bigg)^2
\end{align*}
If we let the potential difference to be expressed in terms of the charge $Q$
as in equation (36)
\begin{align*}
    \Delta\Phi = \frac{4\pi Ql}{A\epsilon}
\end{align*}
Then we get that
\begin{align*}
    F &= -\frac{\epsilon A}{8\pi}\bigg(\frac{\Delta \Phi}{l}\bigg)^2
    = -\frac{\epsilon A}{8\pi}\bigg(\frac{4\pi Q}{A\epsilon}\bigg)^2
    = -\frac{\epsilon A}{8\pi}\frac{16\pi^2 Q^2}{A^2\epsilon^2}
    = -\frac{2\pi Q^2}{A\epsilon}
\end{align*}
\end{proof}

\cleardoublepage
\begin{proof}{\textbf{Exercise 5.}}
Let us take a Gaussian surface shaped as a box where one face is inside the
metallic plate with charge $Q$ and the other face is between the two plates,
then by Gauss law we have that
\begin{align*}
    \int \bm{E}\cdot d\bm{S} &= 4\pi Q_{enc}\\
    \int_{plate} \bm{E}\cdot d\bm{S} + \int_{outside} \bm{E}\cdot d\bm{S} &= 4\pi Q\\
    0 + E A &= 4\pi Q
\end{align*}
Where we used that $E$ inside the metallic plate is 0 and by symmetry we
neglected the integrals with respect to the other faces of the box.
Therefore
\begin{align*}
    E &= \frac{4\pi Q}{A}
\end{align*}
On the other hand, since $E$ is in the direction perpendicular to the plates
we have that
\begin{align*}
    \Delta \Phi = E l = \frac{4\pi Q l}{A}
\end{align*}
Hence the capacitance is 
\begin{align*}
    C = \frac{Q}{\Delta\Phi} = \frac{QA}{4\pi Q l} = \frac{A}{4\pi l}
\end{align*}
\end{proof}
\end{document}